% Standalone build wrapper for mentoring.tex.
% Produces mentoring.pdf directly via pdflatex (never via Preview/Quartz export,
% which corrupts PDF object dictionaries and gets rejected by research.gov).
% Build with: make mentoring   (see Makefile)
\documentclass[11pt,onecolumn]{article}
\usepackage[letterpaper,margin=1in]{geometry}
\usepackage{times}
\usepackage[colorlinks=true,linkcolor=blue,citecolor=blue,urlcolor=blue]{hyperref}
\linespread{1.0}
\setlength{\parindent}{1.5em}
\setlength{\parskip}{0pt}
\pagestyle{empty}
\begin{document}

\section*{Mentoring Plan}
\label{sec:mentoring}

\textbf{Overview and Philosophy.} MD Armanuzzaman is committed to mentoring students in a way that combines rigorous research training, professional development, and inclusive support. His mentoring philosophy is informed by his experience in embedded systems security, trusted execution environments, FPGA-based SoC security, and hands-on cybersecurity training. He aims to prepare students for careers in academia, industry, and public-sector research by helping them develop technical depth, research independence, and strong scientific communication skills.

\textbf{Orientation and Goal Setting.} Each student joining the T\textsubscript{3}SLab begins with an Individual Development Plan (IDP) identifying research interests, skills, milestones, and career goals, mapped onto one of the project's three research thrusts. The IDP is reviewed regularly with the PI to track progress and adjust goals, alongside regular one-on-one mentoring with clear expectations for research conduct, documentation, and collaboration.

\textbf{Research Training and Scholarly Development.} Students will be integrated into the project through research tasks aligned with their background and career goals. The graduate research assistant funded by this project will work across the three research thrusts: constructing application-specific trusted execution domains on FPGA and ARM~CCA platforms (Thrust~1), developing scalable runtime attestation for Cortex-M systems building on the PI's ENOLA toolchain (Thrust~2), and building evidence-guided compartmentalization and recovery mechanisms (Thrust~3). Additional undergraduate and master's students in the lab will contribute to focused technical tasks, prototyping, evaluation, and documentation supporting these thrusts. Mentees will learn how to formulate research questions, review prior work, and write technical reports, and will receive guidance on manuscript preparation for venues such as IEEE S\&P, ACM CCS, USENIX Security, NDSS, IEEE/ACM EMSOFT, ASIACCS, TDSC, and TIFS.

\textbf{Hands-on Research and Professional Communication.} Students will participate in weekly group meetings and project reviews, presenting progress and ideas to develop technical confidence and the ability to explain complex systems to diverse audiences. The PI will coach students on experimental design, manuscript writing, and response to reviewer feedback so they can grow into independent researchers.

\textbf{Career Development and Networking.} Graduate students will gain experience in proposal writing and professional networking. They will contribute to problem formulation, technical planning, and drafting of project materials while receiving career advice on academic, industry, and government pathways. Students will be encouraged to attend workshops, present at conferences, and engage with the project's industry collaborators (below) and other external researchers.

\textbf{Effective Collaboration.} The project benefits from unfunded industrial collaborations that provide complementary expertise: Peter Linder (Octavo Systems) on embedded systems design and FPGA-based SoC platforms, and Malav Vyas (Palo Alto Networks) on SCADA and industrial control system security. Students will have opportunities to engage with these collaborators for technical guidance and exposure to industry perspectives on secure hardware/software co-design and deployment-oriented security research, complementing their day-to-day mentorship within the T\textsubscript{3}SLab.

\textbf{Peer Mentoring and Professional Integrity.} The T\textsubscript{3}SLab currently supports one Ph.D. student, two M.S. students, and two B.S. students, and students will mentor junior researchers within this group as part of a collaborative research culture. Senior students will guide junior students in literature review, implementation, and technical presentation, while all mentees are trained in responsible conduct of research, authorship, and reproducibility. Progress is assessed through IDPs and regular feedback sessions.

\textbf{Expected Outcomes.} This mentoring plan will prepare students to conduct high-quality systems security research, produce scholarly work, and present at conferences. By the end of the project, students are expected to have co-authored publications, contributed to the open-source toolchains developed across the three thrusts, and built the professional foundation for careers in research and engineering.

\end{document}
